% Propietario conservado: /Users/ruben/Documents/New project/output/EXTREMA_MEDIA_RAZON_RECIPROCIDAD_ES_EN_20260918/ARTICULO/ES/source/libro/rectas27.tex
\section{Holonomía finita, configuración de Schläfli y simetría \texorpdfstring{\(E_6\)}{E6}}
\label{x_reciprocidad:p12-83-c20-s6:chap:holonomia-schlafli-e6-nuevo}
\index{holonomía finita}
\index{Schläfli}
\index{E6@\(E_6\)}

La rama de Witt no agota las estructuras excepcionales accesibles desde
APP. La rama paralela de este capítulo parte de las dos polarizaciones
mixtas suma--producto y producto--suma, junto con una orientación de
plaqueta y una sección central declaradas. La curvatura calculada determina
la forma alternante; ésta define la extensión central. No se infiere una
extensión a partir de la sola cardinalidad 27.

\subsection{Curvaturas mixtas de APP}

Sobre \(T_9=(\mathbb Z/9\mathbb Z)^2\), con
\(S(x,y)=x+y\) y \(P(x,y)=xy\), sean
\[
 \Delta_xf=f(x+1,y)-f(x,y),\qquad
 \Delta_yf=f(x,y+1)-f(x,y).
\]
Las conexiones mixtas son
\[
 A^{SP}=(\Delta_xS,\Delta_yP)=(1,x),\qquad
 A^{PS}=(\Delta_xP,\Delta_yS)=(y,1).
\]
Su curvatura discreta es
\[
 F_A=\Delta_xA_y-\Delta_yA_x.
\]

\begin{proposition}[Curvaturas orientadas opuestas]
\label{x_reciprocidad:p12-83-c20-s6:prop:curvaturas-app-e6-nuevo}
\[
 F_{A^{SP}}=1,\qquad F_{A^{PS}}=-1\pmod9.
\]
La circulación sobre un rectángulo orientado de lados \(r,s\) es,
respectivamente, \(rs\) y \(-rs\).
\end{proposition}

\begin{proof}
\(\Delta_xx-\Delta_y1=1\) y
\(\Delta_x1-\Delta_yy=-1\). Un rectángulo contiene \(rs\) plaquetas; al
sumar, las aristas interiores se cancelan y queda la circulación del borde.
\end{proof}

Para desplazamientos \(u=(a,b)\), \(v=(c,d)\), la clase alternante es
\begin{equation}
 \sigma(u,v)=ad-bc\pmod9.
 \label{x_reciprocidad:p12-83-c20-s6:eq:forma-alternante-e6-nueva}
\end{equation}
La orientación fija el signo de \(\sigma\). La sección central empleada
abajo pertenece al calibre de esta realización; cambiarla por un coborde
produce una extensión isomorfa, no un nuevo invariante.

\subsection{Extensión central y reducción ternaria}

Definimos el grupo de Heisenberg nonádico
\begin{equation}
 \mathcal H_9=\{(a,b,z):a,b,z\in\mathbb Z/9\mathbb Z\},
 \quad
 (a,b,z)(c,d,w)=(a+c,b+d,z+w+bc).
 \label{x_reciprocidad:p12-83-c20-s6:eq:heisenberg-9-nuevo}
\end{equation}
El conmutador de dos desplazamientos horizontales es
\[
 [(a,b,0),(c,d,0)]=(0,0,bc-ad),
\]
que registra \(-\sigma(u,v)\). Reduciendo cada coordenada módulo tres se
obtiene
\begin{equation}
 H_3=\{(a,b,z):a,b,z\in\mathbb F_3\},\qquad |H_3|=27,
 \label{x_reciprocidad:p12-83-c20-s6:eq:heisenberg-3-nuevo}
\end{equation}
con la misma ley. El núcleo de
\(\mathcal H_9\to H_3\) es
\((3\mathbb Z/9\mathbb Z)^3\), de cardinal 27.

\subsection{Representación de Weyl--Heisenberg y fase central}

La extensión nonádica admite una realización unitaria explícita. Sea
\[
 \zeta_9:=\exp(2\pi i/9),
 \qquad
 \mathscr H_9^{\rm Sch}:=\mathbb C^9,
\]
con base \(\{|n\rangle:n\in\mathbb Z/9\mathbb Z\}\). Definimos
\begin{equation}
 X|n\rangle:=|n+1\rangle,
 \qquad
 Z|n\rangle:=\zeta_9^n|n\rangle.
 \label{x_reciprocidad:p12-83-c20-s6:eq:weyl-desplazamiento-fase-u029}
\end{equation}
Se tiene
\[
 ZX=\zeta_9XZ
\]
y, por multiplicación,
\begin{equation}
 (X^aZ^b)(X^cZ^d)
 =\zeta_9^{bc}X^{a+c}Z^{b+d}.
 \label{x_reciprocidad:p12-83-c20-s6:eq:cociclo-weyl-nonadico-u029}
\end{equation}

\begin{theorem}[Representación fiel de Weyl--Heisenberg]
\label{x_reciprocidad:p12-83-c20-s6:thm:representacion-weyl-heisenberg-u029}
La aplicación
\begin{equation}
 \rho_{\rm WH}:\mathcal H_9\longrightarrow
 U(\mathscr H_9^{\rm Sch}),
 \qquad
 \rho_{\rm WH}(a,b,z):=\zeta_9^zX^aZ^b,
 \label{x_reciprocidad:p12-83-c20-s6:eq:representacion-weyl-heisenberg-u029}
\end{equation}
es una representación unitaria fiel. Para
\(u=(a,b)\) y \(v=(c,d)\), su conmutador satisface
\begin{equation}
 [\rho_{\rm WH}(u,0),\rho_{\rm WH}(v,0)]
 =\zeta_9^{\,bc-ad}I
 =\zeta_9^{-\sigma(u,v)}I.
 \label{x_reciprocidad:p12-83-c20-s6:eq:conmutador-weyl-holonomia-u029}
\end{equation}
\end{theorem}

\begin{proof}
La identidad \eqref{x_reciprocidad:p12-83-c20-s6:eq:cociclo-weyl-nonadico-u029} reproduce el término
\(bc\) de \eqref{x_reciprocidad:p12-83-c20-s6:eq:heisenberg-9-nuevo}; de ahí la propiedad de
homomorfismo. Si \(\zeta_9^zX^aZ^b=I\), la acción sobre la base obliga
sucesivamente a \(a=0\), \(b=0\) y \(z=0\). Esto prueba fidelidad. Comparar
los productos en los dos órdenes da
\(\zeta_9^{bc-ad}\), que es la fase alternante de
\eqref{x_reciprocidad:p12-83-c20-s6:eq:forma-alternante-e6-nueva}.
\end{proof}

La representación distingue dos observables conjugados. Sean
\[
 Q_{\rm H}:=\operatorname{diag}(0,1,\ldots,8),
 \qquad
 F_{jk}:=\frac13\zeta_9^{jk},
 \qquad
 P_{\rm H}:=F^\dagger Q_{\rm H}F.
\]
Las bases propias de \(Q_{\rm H}\) y \(P_{\rm H}\) son mutuamente
insesgadas, pues \(|F_{jk}|^2=1/9\). Para todo vector unitario
\(\psi\in\mathscr H_9^{\rm Sch}\),
\begin{align}
 \operatorname{Var}_\psi(Q_{\rm H})
 \operatorname{Var}_\psi(P_{\rm H})
 &\ge
 \frac14\left|
 \langle\psi,[Q_{\rm H},P_{\rm H}]\psi\rangle
 \right|^2,
 \label{x_reciprocidad:p12-83-c20-s6:eq:incertidumbre-weyl-u029}\\
 H_{Q_{\rm H}}(\psi)+H_{P_{\rm H}}(\psi)
 &\ge\log9.
 \label{x_reciprocidad:p12-83-c20-s6:eq:incertidumbre-entropica-weyl-u029}
\end{align}
Aquí \(H_{Q_{\rm H}}\) y \(H_{P_{\rm H}}\) son las entropías de Shannon
de las distribuciones de Born en las respectivas bases propias,
con logaritmo natural y convención \(0\log0=0\).
Estos operadores son adimensionales. Una carta física de unidades es una
realización posterior y no interviene en la extensión central.

\subsection{Acción de 10 desplazamientos sobre el grafo de Cayley}

Para un funcional lineal
\(\ell(a,b)=pa+qb\), ponemos
\[
 D_\ell(a,b)=\left(a,b,\frac{ab}{2}+\ell(a,b)\right),
 \qquad Z=(0,0,1),
\]
donde \(2^{-1}=2\) en \(\mathbb F_3\). El conjunto inverso-cerrado
\begin{equation}
 S_\ell=
 \{D_\ell(v):v\in\mathbb F_3^2\setminus\{0\}\}
 \cup\{Z,Z^{-1}\}
 \label{x_reciprocidad:p12-83-c20-s6:eq:conjunto-conexion-e6}
\end{equation}
tiene diez elementos y no contiene la identidad.

\begin{theorem}[Independencia del calibre lineal]
\label{x_reciprocidad:p12-83-c20-s6:thm:independencia-calibre-e6}
Los nueve funcionales \(\ell\) producen grafos de Cayley isomorfos. Los
nueve conjuntos \(S_\ell\) forman una sola órbita bajo
\(\operatorname{Aut}(H_3)\).
\end{theorem}

\begin{proof}
La aplicación
\[
 \phi_\ell(a,b,z)=(a,b,z+\ell(a,b))
\]
es un automorfismo central y satisface
\(\phi_\ell(S_0)=S_\ell\). Para escribir la familia completa en las
coordenadas polarizadas de \cref{x_reciprocidad:p12-83-c20-s6:eq:heisenberg-9-nuevo}, sea
\(c(v,w)=b c\) para \(v=(a,b)\), \(w=(c,d)\). Dado
\(M\in\operatorname{GL}(2,3)\), el defecto
\[
 B_M(v,w)=c(Mv,Mw)-\det(M)c(v,w)
\]
es simétrico, porque las partes alternantes de los dos términos coinciden.
Definimos la corrección cuadrática
\begin{equation}
 q_M(v)=2B_M(v,v),
 \qquad
 q_M(v+w)-q_M(v)-q_M(w)=B_M(v,w),
 \label{x_reciprocidad:p12-83-c20-s6:eq:correccion-cuadratica-automorfismos-h3}
\end{equation}
donde \(2=2^{-1}\) en \(\mathbb F_3\). La familia completa de
automorfismos se escribe entonces
\begin{equation}
 (v,t)\longmapsto
 \bigl(Mv,\det(M)t+q_M(v)+\ell(v)\bigr),
 \qquad M\in\operatorname{GL}(2,3),\
 \ell\in(\mathbb F_3^2)^*.
 \label{x_reciprocidad:p12-83-c20-s6:eq:automorfismos-h3-completos}
\end{equation}
La identidad polar de \cref{x_reciprocidad:p12-83-c20-s6:eq:correccion-cuadratica-automorfismos-h3}
prueba directamente que el producto central se conserva. La familia tiene
\(48\cdot9=432\) elementos. La enumeración exhaustiva de los
\(\binom{13}{5}=1287\) subconjuntos inverso-cerrados de cardinal diez
encuentra exactamente esos nueve con los parámetros que se derivan a
continuación.
\end{proof}

Escribamos \(\underline S=\sum_{s\in S}s\) en
\(\mathbb Z[H_3]\). El conteo de productos ordenados da
\begin{equation}
 \underline S^{\,2}
 =10e+\underline S+5(\underline{H_3}-e-\underline S)
 =5\underline{H_3}+5e-4\underline S.
 \label{x_reciprocidad:p12-83-c20-s6:eq:identidad-grupo-schlafli-nueva}
\end{equation}
En efecto, la identidad recibe diez factorizaciones; un elemento de
\(S\) recibe una; cada uno de los otros dieciséis recibe cinco.

\begin{theorem}[Grafo de las veintisiete rectas]
\label{x_reciprocidad:p12-83-c20-s6:thm:grafo-27-rectas-nuevo}
La matriz de adyacencia \(A_\ell\) de
\(\Gamma_\ell=\operatorname{Cay}(H_3,S_\ell)\) satisface
\begin{equation}
 A_\ell^2=5I-4A_\ell+5J.
 \label{x_reciprocidad:p12-83-c20-s6:eq:identidad-schlafli-nueva}
\end{equation}
Por tanto,
\[
 \Gamma_\ell=\operatorname{srg}(27,10,1,5),
 \qquad
 \operatorname{Spec}(A_\ell)=\{10^1,1^{20},(-5)^6\}.
\]
Es isomorfo al grafo de intersección de las veintisiete rectas de una
superficie cúbica lisa.
\end{theorem}

\begin{proof}
La representación regular de
\cref{x_reciprocidad:p12-83-c20-s6:eq:identidad-grupo-schlafli-nueva} produce
\cref{x_reciprocidad:p12-83-c20-s6:eq:identidad-schlafli-nueva}. Sus entradas dicen que todo vértice
tiene diez vecinos, que un par adyacente tiene un vecino común y un par no
adyacente tiene cinco. Sobre la recta constante, \(A_\ell\) actúa por 10;
en su complemento, \(J=0\) y
\((A_\ell-I)(A_\ell+5I)=0\). Traza cero fija las multiplicidades 20 y 6.

En el modelo de una cúbica obtenida al soplar seis puntos de
\(\mathbb P^2\), las 27 clases son
\[
 E_i,\qquad L_{ij}=H-E_i-E_j,\qquad
 Q_i=2H-\sum_{j\ne i}E_j.
\]
Con \(H^2=1\), \(E_i^2=-1\), \(H\cdot E_i=0\) y
\(E_i\cdot E_j=0\) para \(i\ne j\), la matriz de adyacencia definida
por intersección entre clases distintas tiene los mismos parámetros.
La matriz completa de intersección tiene diagonal \(-1\).
La identificación con la configuración clásica usa el
teorema de unicidad de esta realización de las 27 rectas
\cite{x_reciprocidad:Manivel2005}, con el dominio y la acción allí especificados.
\end{proof}

La afirmación de isomorfía puede hacerse completamente efectiva mediante
una biyección calculada. El etiquetado no es absoluto: otra sección lineal
o una automorfía de la superficie produce otra tabla de la misma clase.
\begin{longtable}{@{}c@{\,}c@{\,}c@{\qquad}l@{}}
\caption{Bijección explícita entre \(H_3\) y las veintisiete clases de
rectas.}
\label{x_reciprocidad:p12-83-c20-s6:tab:h3-27-rectas-u029}\\
\toprule
\(a\)&\(b\)&\(z\)&clase de recta\\
\midrule
\endfirsthead
\multicolumn{4}{@{}l}{\small\itshape Cuadro \thetable\ (continuación)}\\
\toprule
\(a\)&\(b\)&\(z\)&clase de recta\\
\midrule
\endhead
0&0&0&\(E_1\)\\
0&0&1&\(L_{12}\)\\
0&0&2&\(Q_2\)\\
0&1&0&\(L_{13}\)\\
0&1&1&\(Q_1\)\\
0&1&2&\(E_3\)\\
0&2&0&\(Q_3\)\\
0&2&1&\(E_2\)\\
0&2&2&\(L_{23}\)\\
1&0&0&\(L_{14}\)\\
1&0&1&\(L_{35}\)\\
1&0&2&\(L_{26}\)\\
1&1&0&\(L_{36}\)\\
1&1&1&\(L_{24}\)\\
1&1&2&\(L_{15}\)\\
1&2&0&\(L_{25}\)\\
1&2&1&\(L_{16}\)\\
1&2&2&\(L_{34}\)\\
2&0&0&\(Q_4\)\\
2&0&1&\(L_{46}\)\\
2&0&2&\(E_6\)\\
2&1&0&\(E_5\)\\
2&1&1&\(Q_6\)\\
2&1&2&\(L_{56}\)\\
2&2&0&\(L_{45}\)\\
2&2&1&\(E_4\)\\
2&2&2&\(Q_5\)\\
\bottomrule
\end{longtable}

Como control local, los diez vecinos de \(e=(0,0,0)\) son los elementos de
\(S_0\). La tabla los transforma en las diez clases que cortan a \(E_1\):
las cinco \(L_{1j}\) y las cinco \(Q_j\), con \(2\le j\le6\). La
verificación completa compara las 729 entradas de ambas matrices de
adyacencia.

El complemento \(B_\ell=J-I-A_\ell\) satisface
\begin{equation}
 B_\ell^2=8I+2B_\ell+8J,
 \qquad
 \operatorname{srg}(27,16,10,8),
 \qquad
 \operatorname{Spec}(B_\ell)=\{16^1,4^6,(-2)^{20}\}.
 \label{x_reciprocidad:p12-83-c20-s6:eq:complemento-schlafli}
\end{equation}
Éste es el grafo de Schläfli en la convención de adyacencia complementaria;
\(A_\ell\) es el grafo de intersección de parámetros
\((27,10,1,5)\).

Cada arista pertenece a un triángulo único, pues \(\lambda=1\). El número
de triángulos es
\[
 \frac{27\cdot10}{6}=45.
\]
El grupo completo de automorfismos tiene orden
\[
 27\cdot1920=51840
\]
y, por su acción sobre la configuración de rectas, se identifica con
\(W(E_6)\), no sólo por igualdad de órdenes \cite{x_reciprocidad:Manivel2005}.

\subsection{El retículo de raíces que realiza \texorpdfstring{\(E_6\)}{E6}}

Para hacer explícito el objeto sobre el que actúa ese grupo, sea
\begin{equation}
 \operatorname{Pic}(X)
 =\mathbb ZH\oplus\bigoplus_{i=1}^6\mathbb ZE_i,
 \qquad
 H^2=1,\qquad E_i^2=-1,\qquad H\cdot E_i=E_i\cdot E_j=0\ (i\ne j),
 \label{x_reciprocidad:p12-83-c20-s6:eq:picard-cubica-e6}
\end{equation}
el retículo de Picard de la explosión de seis puntos generales de
\(\mathbb P^2\). Su clase canónica es
\[
 K_X=-3H+E_1+\cdots+E_6.
\]
El complemento ortogonal
\begin{equation}
 L_{E_6}=K_X^\perp\subset\operatorname{Pic}(X),
 \qquad (x,y)_{E_6}:=-x\cdot y,
 \label{x_reciprocidad:p12-83-c20-s6:eq:reticulo-raices-e6}
\end{equation}
es positivo definido de rango seis. Una base de raíces simples es
\begin{equation}
 \begin{aligned}
 \alpha_1&=E_1-E_2,& \alpha_2&=E_2-E_3,&
 \alpha_3&=E_3-E_4,\\
 \alpha_4&=E_4-E_5,& \alpha_5&=E_5-E_6,&
 \alpha_6&=H-E_1-E_2-E_3.
 \end{aligned}
 \label{x_reciprocidad:p12-83-c20-s6:eq:raices-simples-e6}
\end{equation}
Cada \(\alpha_i\) es ortogonal a \(K_X\), tiene norma dos y sus productos
son \(-1\) exactamente para los pares adyacentes del diagrama
\(E_6\): la cadena \(1-2-3-4-5\), con el nodo \(6\) unido al nodo \(3\).
Por tanto, su matriz de Gram es la matriz de Cartan de \(E_6\).

\begin{proposition}[Las setenta y dos raíces y sus reflexiones]
\label{x_reciprocidad:p12-83-c20-s6:prop:setenta-dos-raices-e6}
Las raíces de \(L_{E_6}\) son
\begin{equation}
 \begin{aligned}
 &E_i-E_j &&(i\ne j),\\
 &\pm(H-E_i-E_j-E_k) &&(1\le i<j<k\le6),\\
 &\pm(2H-E_1-\cdots-E_6),
 \end{aligned}
 \label{x_reciprocidad:p12-83-c20-s6:eq:enumeracion-raices-e6}
\end{equation}
en total \(30+40+2=72\). Para cada raíz \(\alpha\),
\begin{equation}
 s_\alpha(x)=x+(x\cdot\alpha)\alpha
 \label{x_reciprocidad:p12-83-c20-s6:eq:reflexion-raiz-e6}
\end{equation}
preserva la intersección y \(K_X\). Las seis reflexiones simples generan
\(W(E_6)\) y permutan fielmente las 27 clases
\(E_i,L_{ij},Q_i\) de \cref{x_reciprocidad:p12-83-c20-s6:thm:grafo-27-rectas-nuevo}.
\end{proposition}

\begin{proof}
Sustituir cada clase de \cref{x_reciprocidad:p12-83-c20-s6:eq:enumeracion-raices-e6} en
\cref{x_reciprocidad:p12-83-c20-s6:eq:picard-cubica-e6} da \(K_X\cdot\alpha=0\) y
\(\alpha^2=-2\). La cuenta de clases es la indicada. La fórmula de reflexión
es la reflexión ortogonal para la forma \(-(\cdot)\); preserva \(K_X\) porque
\(K_X\cdot\alpha=0\). El sistema simple de
\cref{x_reciprocidad:p12-83-c20-s6:eq:raices-simples-e6} genera todas las raíces y su grupo de Coxeter es
\(W(E_6)\). Su acción sobre las 27 clases de curvas excepcionales conserva la
intersección, de modo que coincide con la acción sobre el grafo ya
construido \cite{x_reciprocidad:Manivel2005}.
\end{proof}

\subsection{Estratificación de las 729 palabras}

Escribimos una palabra de seis trits como
\[
 w=(\ell_1,h_1,\ell_2,h_2,\ell_3,h_3),
\]
y definimos dos elementos de \(H_3\),
\[
 L(w)=(\ell_1,\ell_2,\ell_3),\qquad
 H(w)=(h_1,h_2,h_3),
\]
usando la ley central en sus productos. Según
\(L(w)^{-1}H(w)\), las 729 palabras se dividen en
\begin{align*}
 \mathcal D_\ell&=\{w:L(w)^{-1}H(w)=e\},\\
 \mathcal I_\ell&=\{w:L(w)^{-1}H(w)\in S_\ell\},\\
 \mathcal S_\ell&=\{w:L(w)^{-1}H(w)\notin\{e\}\cup S_\ell\}.
\end{align*}

\begin{theorem}[Partición Hensel--Schläfli]
\label{x_reciprocidad:p12-83-c20-s6:thm:particion-hensel-schlafli-nueva}
\begin{equation}
 729=27+270+432,
 \label{x_reciprocidad:p12-83-c20-s6:eq:particion-729-e6-nueva}
\end{equation}
correspondientes a diagonal, incidencia y alabeamiento.
\end{theorem}

\begin{proof}
Para cada uno de los 27 valores de \(L(w)\), hay un único valor de
\(H(w)\) con diferencia identidad, diez con diferencia en \(S_\ell\) y
dieciséis restantes. Se multiplica \(27\) por \(1,10,16\).
\end{proof}

El cardinal \(432\) del estrato alabeado coincide numéricamente con la
multiplicidad de la región transversal de \(\pi\) en
\cref{x_reciprocidad:pi-estrella:figura}. No se identifica por ello ambos objetos: aquí
se cuentan palabras clasificadas por \(L(w)^{-1}H(w)\), mientras allí se
cuentan preimágenes de un registro \(U_6\). Sin un entrelazador explícito y
equivariante entre esos dominios, \(432=432\) es sólo un control cardinal.

\subsection{Necesidad de la fase central}

La enumeración de todos los subconjuntos inverso-cerrados de cardinal diez
en los cinco grupos de orden 27 da
\[
\begin{array}{c|c|c}
\text{grupo}&\text{abeliano}&
\#\operatorname{srg}(27,10,1,5)\\ \hline
C_{27}&\text{sí}&0\\
C_9\times C_3&\text{sí}&0\\
C_3^3&\text{sí}&0\\
C_9\rtimes C_3&\text{no}&9\\
H_3&\text{no}&9.
\end{array}
\]
Así, conservar 27 etiquetas pero abelianizar la composición destruye la
configuración. No bastan la cardinalidad ni el grado.

Si \(\mathcal T\) es el conjunto de 45 triángulos, el cúbico
\[
 \mathcal C_\Gamma(x_1,\ldots,x_{27})
 =\sum_{\{i,j,k\}\in\mathcal T}x_ix_jx_k
\]
recupera la adyacencia, porque cada arista está en un triángulo único. Su
estabilizador permutacional es exactamente
\[
 \operatorname{Aut}(\Gamma_\ell)\cong W(E_6).
\]

Para el diagrama siguiente se escriben
\(\Gamma_{27}=\Gamma_\ell\) y \(\mathcal T_{45}=\mathcal T\).
\begin{figure}[htbp]
\centering
\[
 \begin{gathered}
 \text{curvaturas APP }\pm1
 \longrightarrow\sigma
 \longrightarrow\mathcal H_9
 \longrightarrow H_3,\\[3pt]
 H_3
 \longrightarrow\Gamma_{27}
 \longrightarrow\mathcal T_{45}
 \longrightarrow W(E_6).
 \end{gathered}
\]
\caption{Cadena causal de la rama de holonomía mixta. Es paralela a
\(C_W\to W_{12}\to M_{12}\) y no un eslabón de ella.}
\label{x_reciprocidad:p12-83-c20-s6:fig:cadena-holonomia-e6-u029}
\end{figure}

\begin{criterion}[Criterios de refutación de la rama \(E_6\)]
\label{x_reciprocidad:p12-83-c20-s6:crit:refutacion-rama-e6-u029}
La rama queda refutada si las curvaturas no son opuestas, si el cociclo
central se borra, si la corrección cuadrática
\eqref{x_reciprocidad:p12-83-c20-s6:eq:correccion-cuadratica-automorfismos-h3} no conserva el producto,
si \eqref{x_reciprocidad:p12-83-c20-s6:eq:identidad-grupo-schlafli-nueva} no produce los coeficientes
\((10,1,5)\), si un grupo abeliano produce la misma configuración, si la
partición no suma 729, si la tabla
\ref{x_reciprocidad:p12-83-c20-s6:tab:h3-27-rectas-u029} no entrelaza las matrices de adyacencia o si
\(W(E_6)\) se identifica sólo por su orden y no por su acción sobre las
veintisiete clases y las 72 raíces.
\end{criterion}
